% ============================================================================
% BAB I - PENDAHULUAN
% ============================================================================
\chapter{PENDAHULUAN}

\section{Latar Belakang}

Perkembangan teknologi informasi yang pesat mendorong kebutuhan akan
sistem yang mampu menangani volume data dan lalu lintas query yang
tinggi secara efisien. Pada aplikasi \textit{e-commerce} dan sistem
pengelolaan data seperti Komposin (sebuah platform bank sampah dan
\textit{e-commerce} pengelolaan sampah), beban operasi \textit{read}
(pembacaan data) umumnya jauh lebih tinggi dibandingkan operasi
\textit{write} (penulisan data). Pengguna lebih sering melakukan
pencarian produk, melihat katalog, membaca artikel, dan mengecek
riwayat transaksi dibandingkan melakukan pembelian atau pendaftaran
akun. Karakteristik beban kerja seperti ini lazim ditemukan pada
hampir semua aplikasi web modern, di mana rasio \textit{read} terhadap
\textit{write} dapat mencapai 80:20 atau bahkan 90:10. Oleh karena
itu, strategi pengelolaan database perlu disesuaikan agar sumber daya
komputasi dialokasikan secara proporsional terhadap jenis beban yang
dominan.

Apabila hanya menggunakan satu \textit{server} database (PostgreSQL
\textit{single instance}), seluruh beban query, baik \textit{read}
maupun \textit{write}, diproses oleh \textit{server} yang sama.
Kondisi ini menimbulkan beberapa permasalahan, antara lain:

\begin{enumerate}
    \item \textbf{Bottleneck performa.} Ketika jumlah pengguna
    meningkat, \textit{server} tunggal tidak mampu menangani lonjakan
    query secara bersamaan, menyebabkan waktu respons yang lambat.
    Pengguna akan merasakan halaman yang lama dimuat, terutama pada
    jam-jam sibuk ketika banyak pengguna mengakses katalog dan
    melakukan transaksi secara bersamaan. Kondisi ini dapat menurunkan
    kepuasan pengguna dan berpotensi mengurangi pendapatan aplikasi.
    \item \textbf{Tidak ada toleransi kesalahan (\textit{fault
    tolerance}).} Jika \textit{server} database mengalami kegagalan
    (\textit{down}), seluruh aplikasi berhenti beroperasi karena tidak
    ada cadangan. Downtime yang terjadi dapat berlangsung dalam hitungan
    menit hingga jam, tergantung pada kecepatan pemulihan server.
    \item \textbf{Beban \textit{read} mengganggu operasi
    \textit{write}.} Query \texttt{SELECT} yang berat (misalnya laporan
    transaksi) dapat memperlambat query \texttt{INSERT/UPDATE} yang
    membutuhkan komitmen data segera. Akibatnya, proses pembelian atau
    pendaftaran pengguna baru dapat tertunda oleh query pembacaan yang
    seharusnya tidak mendesak.
\end{enumerate}

Untuk mengatasi permasalahan tersebut, diterapkan arsitektur
\textbf{\textit{read-write splitting}} dengan \textbf{replikasi
database}. Mekanisme ini memisahkan query \textit{write}
(\texttt{INSERT}, \texttt{UPDATE}, \texttt{DELETE}) yang dikirim ke
\textbf{PostgreSQL Master}, dan query \textit{read} (\texttt{SELECT})
yang didistribusikan ke satu atau lebih \textbf{PostgreSQL Replica}.
Sebagai \textit{middleware} yang mengatur distribusi query, digunakan
\textbf{Pgpool-II}, sebuah \textit{connection pooler} dan
\textit{load balancer} yang secara \textit{transparan} mengarahkan
query ke \textit{backend} yang tepat berdasarkan jenis operasinya.
Pgpool-II juga menyediakan mekanisme \textit{health check} dan
\textit{failover} otomatis, sehingga jika salah satu Replica gagal,
traffic akan dialihkan ke Replica lain tanpa mengganggu aplikasi.
Keseluruhan arsitektur di-\textit{deploy} menggunakan Docker Compose
untuk memastikan konsistensi lingkungan dan kemudahan replikasi
konfigurasi.

Dengan arsitektur ini, aplikasi Komposin diharapkan mampu menangani
beban \textit{read} yang tinggi tanpa mengorbankan performa
\textit{write}. Selain itu, sistem diharapkan tetap beroperasi
meskipun salah satu \textit{replica} mengalami kegagalan
(\textit{high availability}), karena Pgpool-II akan secara otomatis
mengalihkan koneksi ke node yang masih sehat. Arsitektur ini juga
memungkinkan distribusi beban query secara merata ke beberapa
\textit{node} database melalui mekanisme \textit{load balancing}.
Keempat, penggunaan \textit{connection pooling} mengurangi overhead
pembukaan koneksi baru yang berulang. Kelima, pemisahan beban
\textit{read} dan \textit{write} memungkinkan masing-masing server
dioptimalkan secara independen sesuai perannya.

\section{Ruang Lingkup Sistem}

\subsection{Aplikasi yang Dibangun}

Aplikasi Komposin adalah platform berbasis web yang dikembangkan
menggunakan \textbf{Laravel 10} (PHP) dengan database
\textbf{PostgreSQL 16}. Aplikasi ini memiliki dua peran pengguna:
\textbf{Admin} dan \textbf{Pelanggan}, di mana Admin mengelola data
produk, artikel, reward, transaksi, rute pengambilan sampah, serta
laporan. Sementara itu, Pelanggan dapat melihat katalog produk,
melakukan pembelian, menukarkan reward, melihat riwayat transaksi, dan
berlangganan layanan pengambilan sampah. Seluruh fitur tersebut
mengandalkan operasi database yang intensif, baik \textit{read}
maupun \textit{write}, sehingga menjadi kandidat ideal untuk
penerapan \textit{read-write splitting}. Aplikasi dikemas dalam
container Docker yang berjalan di atas Nginx dan Supervisor untuk
memastikan keandalan layanan.

\subsection{Fitur Utama Aplikasi}

\begin{itemize}
    \item \textbf{Manajemen Katalog Produk.} Admin dapat menambah,
    mengubah, dan menghapus produk. Pelanggan dapat melihat daftar
    produk. Fitur ini didominasi oleh operasi \texttt{SELECT} yang
    tinggi, terutama saat banyak pengguna menjelajahi katalog secara
    bersamaan.
    \item \textbf{Manajemen Transaksi.} Pelanggan melakukan pembelian
    produk melalui keranjang belanja dan \textit{check-out}. Admin
    mengelola status transaksi. Setiap transaksi melibatkan kombinasi
    query \texttt{INSERT} (membuat pesanan) dan \texttt{SELECT}
    (memvalidasi stok dan data pelanggan).
    \item \textbf{Manajemen Reward.} Pelanggan menukarkan poin dengan
    reward, dan Admin mengonfirmasi penukaran reward. Proses ini
    memerlukan validasi saldo poin melalui \texttt{SELECT} dan
    pencatatan penukaran melalui \texttt{INSERT}.
    \item \textbf{Manajemen Artikel.} Admin membuat dan mengelola
    artikel edukasi seputar pengelolaan sampah. Pelanggan membaca
    artikel yang mayoritas hanya memerlukan operasi \texttt{SELECT}.
    \item \textbf{Berlangganan.} Pelanggan dapat berlangganan layanan
    pengambilan sampah berkala. Sistem mencatat jadwal dan status
    berlangganan melalui operasi \texttt{INSERT} dan \texttt{UPDATE}.
    \item \textbf{Pelaporan.} Admin melihat laporan transaksi dan
    penghasilan dari layanan berlangganan. Laporan melibatkan query
    \texttt{SELECT} agregat yang kompleks dan berpotensi membebani
    database.
    \item \textbf{Pengambilan Sampah.} Pelanggan menjadwalkan
    pengambilan sampah dan mendapatkan poin sebagai imbalan. Fitur ini
    melibatkan pencatatan jadwal (\texttt{INSERT}) dan pembaruan poin
    (\texttt{UPDATE}).
\end{itemize}

\subsection{Batasan Implementasi Sistem}

Batasan implementasi pada final project ini ditetapkan untuk menjaga
fokus dan kelayakan teknis dalam lingkup mata kuliah Basis Data
Terdistribusi. Sistem menggunakan \textbf{1 (satu) PostgreSQL Master}
sebagai \textit{node} utama yang menerima seluruh operasi
\textit{write}, serta \textbf{2 (dua) PostgreSQL Replica} sebagai
\textit{node} pembaca yang menerima operasi \textit{read}. Satu
\textbf{Pgpool-II} berperan sebagai \textit{middleware} pengatur
\textit{connection pooling}, \textit{load balancing}, dan
\textit{read-write splitting}. Seluruh layanan di-\textit{deploy}
menggunakan \textbf{Docker Compose} dalam satu \textit{host} dengan
jaringan \textit{bridge} internal, sehingga konfigurasi dapat
direproduksi dengan mudah di lingkungan lain. Replikasi menggunakan
metode \textbf{\textit{streaming replication}} secara
\textbf{asinkron} yang sudah tersedia secara \textit{native} pada
PostgreSQL 16. Aplikasi \textbf{tidak} terhubung langsung ke
PostgreSQL Master maupun Replica. Seluruh koneksi database
diarahkan ke Pgpool-II sebagai \textit{single entry point}.
Pendekatan ini memastikan bahwa mekanisme \textit{read-write
splitting} berjalan secara transparan tanpa modifikasi kode aplikasi.
